\[
A=\begin{pmatrix}
1&0&0&1&-1&1\\0&1&-1&-1&1&0\\0&1&0&-1&0&0\\
0&2&-1&-1&0&1\\0&2&0&-1&0&0\\0&1&1&-1&0&0
\end{pmatrix},\quad
B=\begin{pmatrix}
1&0&0&0&0&0\\0&0&-1&0&1&0\\0&0&-1&0&0&1\\
0&0&-1&0&0&0\\0&1&-1&0&0&0\\0&0&-1&1&0&0
\end{pmatrix}.
\]
Put $M_1=A^t$, $M_2=B^t$, $M_0=(M_1M_2)^{-1}$ and
\[
 v_1=(1,3,-2,-1,2,0)^t,\qquad v_2=(-1,0,0,0,0,0)^t.
\]
Theorem~\ref{night:P1.1} gives their orders, the primitive cusp
marking, the coinvariant homomorphism $\ell(x)=x_0$, and freeness of
every nonidentity power of the specified affine actions. Consequently
\[
 (T^6\times\Delta_u)/
 \langle(x,u)\mapsto(M_ix+v_i/m_i,\zeta_{m_i}u)\rangle,
 \qquad (m_1,m_2)=(5,4),
\]
are smooth oriented finite caps. On the universal angular cover
$u=r e^{2\pi ia}$ their actual marked collar maps are
\[
 \alpha_i(x,u)=x-av_i,
 \qquad
 \alpha_i(M_ix+v_i/m_i,\zeta_{m_i}u)=M_i\alpha_i(x,u).
\]
At the cusp we take the actual index-one Kuhn cap of
Theorem~\ref{night:P1.9}, in the marking of
Proposition~\ref{maps:src:F.2.1}. The three displayed boundary bundles
therefore glue to the marked torus bundle over the intervening pair
of pants, with zero additional clutch translation. Every boundary
is paired. The result is a closed connected smooth eight-manifold,
oriented by the cusp complex structure and the continued fibre/base
orientation.
